\begin{abstract}
This paper presents a discrete model of reality built entirely from experience. Reality is modeled as a growing population of \emph{vibes}, indivisible units of experience, related only by the \emph{note} (the primitive relation of one vibe experiencing another), each carrying a finite \emph{tone}, all updating in discrete steps called \emph{beats}. Two constraints govern the whole construction. The first, discreteness, requires that the base be finite, with the continuum permitted only as an emergent large-scale approximation. The second requires that the vibe be the only kind of thing, so that there is no separate link, field, force, or particle substrate, and the only relation is the note. From these two constraints and a single axiom, that to exist is to experience or be experienced, the model derives, in order, its local update dynamics, the energy that experience minimizes, the discrete time it runs in, the necessity of its own beginning, the unbounded growth of experience that follows, the emergence of selves as self-maintaining coherent patterns, the emergence of space and its hyperbolic geometry, the mechanics of how an agent moves through a substrate woven from other agents, and a route toward building matter without adding any new primitive. Physical law is recovered as a large-scale limit. Experience is not recovered. It is the foundation. This work is a philosophical and mathematical model with a concrete, simulable build target rather than an empirically tested theory.
\medskip

\noindent\textbf{Status and caveats.} This is exploratory work, a feasibility study rather than a validated result. The aim is to see whether a model of this kind can be made to hang together at all, and so far it looks promising. Much validation and further work remain. Several of the experiments reported here are preliminary, and some may turn out to be inaccurate or improperly set up. Nothing here should be read as established or taken too rigorously at this stage. It is offered in the spirit of curiosity and invitation, a direction that looks worth exploring rather than a finished theory or a claim of discovery.
\end{abstract}
